\documentclass[12pt,a4paper]{article}

% Packages
\usepackage[utf8]{inputenc}
\usepackage[T1]{fontenc}
\usepackage{amsmath}
\usepackage{amssymb}
\usepackage{amsthm}
\usepackage{graphicx}
\usepackage{booktabs}
\usepackage{algorithm}
\usepackage{algpseudocode}
\usepackage{hyperref}
\usepackage{cleveref}
\usepackage{siunitx}
\usepackage{listings}
\usepackage{xcolor}
\usepackage{geometry}
\usepackage{caption}
\usepackage{subcaption}
\usepackage{tikz}
\usepackage{pgfplots}
\pgfplotsset{compat=1.18}

% Page layout
\geometry{
    left=25mm,
    right=25mm,
    top=30mm,
    bottom=30mm
}

% Code listing style
\lstset{
    basicstyle=\ttfamily\small,
    keywordstyle=\color{blue},
    commentstyle=\color{green!60!black},
    stringstyle=\color{red},
    showstringspaces=false,
    breaklines=true,
    numbers=left,
    numberstyle=\tiny\color{gray},
    frame=single,
    language=Python
}

% Theorem environments
\newtheorem{theorem}{Theorem}
\newtheorem{lemma}{Lemma}
\newtheorem{definition}{Definition}
\newtheorem{proposition}{Proposition}

% Title information
\title{\textbf{Hierarchical Control of Cooperative Quadrotor Payload Transport}\\
\Large A 6-DOF Rigid Body Dynamics Simulation with Geometric Attitude Control}

\author{
    % Add your name(s) here
}

\date{\today}

\begin{document}

\maketitle

\begin{abstract}
This report presents a comprehensive study of cooperative payload transport using multiple quadrotor drones with hierarchical control architecture. We implement a high-fidelity 6-DOF rigid body dynamics simulation for a cubic payload suspended by four quadrotors, employing geometric attitude control on SO(3) combined with PID position control. The system features quintic polynomial trajectory generation for smooth motion profiles, first-order actuator dynamics modeling motor lag, and realistic noise injection including drone mass variance, thrust noise, and attachment point errors. Our hierarchical control architecture consists of an outer-loop PID position controller and an inner-loop geometric attitude controller with least-squares thrust allocation. The simulation framework supports multi-phase trajectory planning including takeoff, translation, and landing maneuvers. Experimental results demonstrate robust performance under realistic operating conditions with configurable noise parameters and controller gains.
\end{abstract}

\tableofcontents
\clearpage

% =============================================================================
\section{Introduction}
% =============================================================================

\subsection{Motivation}
Cooperative aerial manipulation has emerged as a critical area of robotics research with applications in construction, disaster response, and logistics. Multi-drone systems can transport payloads exceeding individual drone capacity, offering advantages in payload-to-weight ratio, redundancy, and operational flexibility compared to single large vehicles.

The control of cooperative payload transport presents several challenges:
\begin{itemize}
    \item \textbf{Coupled dynamics}: The payload and drones form a closed kinematic chain with strong dynamic coupling
    \item \textbf{Underactuation}: Quadrotors are underactuated systems with four control inputs governing six degrees of freedom
    \item \textbf{Nonlinear geometry}: Attitude dynamics evolve on the special orthogonal group SO(3), requiring careful treatment to avoid singularities
    \item \textbf{Thrust allocation}: Desired wrench must be distributed among multiple actuators with saturation constraints
    \item \textbf{Model uncertainty}: Real systems exhibit mass variance, actuator dynamics, and attachment errors
\end{itemize}

\subsection{Contributions}
This work makes the following contributions:
\begin{enumerate}
    \item A complete 6-DOF rigid body simulation framework implementing Newton-Euler dynamics with Rodrigues integration for rotational motion
    \item Hierarchical control architecture combining geometric attitude control with PID position regulation
    \item Quintic polynomial trajectory generation ensuring C$^2$ continuity with zero endpoint velocities
    \item Realistic noise modeling including drone mass variance, actuator noise, and geometric uncertainties
    \item Comprehensive simulation environment with configurable parameters and real-time telemetry
\end{enumerate}

\subsection{Organization}
The remainder of this report is organized as follows: \Cref{sec:background} provides mathematical background on rigid body dynamics and Lie group theory. \Cref{sec:modeling} describes the system model including payload dynamics, drone actuators, and disturbances. \Cref{sec:control} presents the hierarchical control architecture. \Cref{sec:trajectory} details trajectory generation. \Cref{sec:implementation} discusses implementation details. \Cref{sec:experiments} presents experimental results. \Cref{sec:conclusion} concludes.

% =============================================================================
\section{Mathematical Background}
\label{sec:background}
% =============================================================================

\subsection{Rigid Body Dynamics}

\subsubsection{Configuration Space}
The configuration of a rigid body is described by its position $\mathbf{p} \in \mathbb{R}^3$ and orientation $\mathbf{R} \in \text{SO}(3)$, where SO(3) is the special orthogonal group:
\begin{equation}
    \text{SO}(3) = \{\mathbf{R} \in \mathbb{R}^{3 \times 3} : \mathbf{R}^T\mathbf{R} = \mathbf{I}, \det(\mathbf{R}) = 1\}
\end{equation}

The tangent space to SO(3) is the Lie algebra $\mathfrak{so}(3)$ of skew-symmetric matrices:
\begin{equation}
    \mathfrak{so}(3) = \{\boldsymbol{\Omega} \in \mathbb{R}^{3 \times 3} : \boldsymbol{\Omega}^T = -\boldsymbol{\Omega}\}
\end{equation}

\subsubsection{Translational Dynamics}
The linear momentum evolves according to Newton's second law:
\begin{equation}
    m\dot{\mathbf{v}} = \sum_{i=1}^{n} \mathbf{F}_i + m\mathbf{g}
    \label{eq:newton}
\end{equation}
where $m$ is total mass, $\mathbf{v}$ is velocity, $\mathbf{F}_i$ are external forces, and $\mathbf{g} = [0, 0, -9.80665]^T$ m/s$^2$ is gravitational acceleration.

\subsubsection{Rotational Dynamics}
Angular momentum evolves according to Euler's equation:
\begin{equation}
    \mathbf{I}\dot{\boldsymbol{\omega}} = \sum_{i=1}^{n} \boldsymbol{\tau}_i - \boldsymbol{\omega} \times (\mathbf{I}\boldsymbol{\omega})
    \label{eq:euler}
\end{equation}
where $\mathbf{I} \in \mathbb{R}^{3 \times 3}$ is the moment of inertia tensor, $\boldsymbol{\omega} \in \mathbb{R}^3$ is angular velocity, and $\boldsymbol{\tau}_i$ are external torques. The term $\boldsymbol{\omega} \times (\mathbf{I}\boldsymbol{\omega})$ represents gyroscopic (Coriolis) effects.

For a uniform cube with mass $m$ and side length $a$:
\begin{equation}
    \mathbf{I} = \frac{ma^2}{6}\mathbf{I}_3
    \label{eq:cube_inertia}
\end{equation}

\subsubsection{Orientation Integration}
The orientation matrix evolves according to:
\begin{equation}
    \dot{\mathbf{R}} = \mathbf{R}[\boldsymbol{\omega}]_\times
\end{equation}
where $[\cdot]_\times : \mathbb{R}^3 \to \mathfrak{so}(3)$ is the skew-symmetric matrix operator:
\begin{equation}
    [\mathbf{x}]_\times = \begin{bmatrix}
        0 & -x_3 & x_2 \\
        x_3 & 0 & -x_1 \\
        -x_2 & x_1 & 0
    \end{bmatrix}
\end{equation}

We integrate using the exponential map:
\begin{equation}
    \mathbf{R}(t + \Delta t) = \mathbf{R}(t) \exp([\boldsymbol{\omega}\Delta t]_\times)
    \label{eq:rotation_integration}
\end{equation}
implemented via the Rodrigues formula in \texttt{scipy.spatial.transform.Rotation}.

\subsection{Special Orthogonal Group SO(3)}

\begin{definition}[Special Orthogonal Group]
The special orthogonal group SO(3) is a Lie group representing proper rotations in three-dimensional space. It is a smooth manifold diffeomorphic to $\mathbb{RP}^3$ with group operation given by matrix multiplication.
\end{definition}

Key properties:
\begin{itemize}
    \item \textbf{Non-commutativity}: $\mathbf{R}_1\mathbf{R}_2 \neq \mathbf{R}_2\mathbf{R}_1$ in general
    \item \textbf{Compactness}: SO(3) is a compact manifold
    \item \textbf{Connectedness}: SO(3) is path-connected
    \item \textbf{Singularity-free}: No coordinate singularities (unlike Euler angles)
\end{itemize}

% =============================================================================
\section{System Modeling}
\label{sec:modeling}
% =============================================================================

\subsection{Payload Dynamics}

The payload is modeled as a rigid cube with:
\begin{itemize}
    \item Mass: $m_c = \SI{10}{kg}$
    \item Side length: $a = \SI{1}{m}$
    \item Moment of inertia: $\mathbf{I} = \frac{10}{6}\mathbf{I}_3 \approx 1.667\mathbf{I}_3$ kg$\cdot$m$^2$
\end{itemize}

Four attachment points are located at the top corners:
\begin{equation}
    \mathbf{r}_i^{\text{body}} = \frac{a}{2}\begin{bmatrix} \pm 1 \\ \pm 1 \\ 1 \end{bmatrix}, \quad i \in \{0,1,2,3\}
    \label{eq:attachment_local}
\end{equation}

World coordinates are obtained via:
\begin{equation}
    \mathbf{r}_i^{\text{world}} = \mathbf{p} + \mathbf{R}\mathbf{r}_i^{\text{body}}
    \label{eq:attachment_world}
\end{equation}

\subsubsection{External Forces and Torques}
Each drone $i$ applies force $\mathbf{F}_i$ at attachment point $\mathbf{r}_i$, producing torque:
\begin{equation}
    \boldsymbol{\tau}_i = (\mathbf{r}_i - \mathbf{p}) \times \mathbf{F}_i
    \label{eq:torque}
\end{equation}

Total external wrench:
\begin{align}
    \mathbf{F}_{\text{total}} &= \sum_{i=0}^{3} \mathbf{F}_i + (m_c + m_d)\mathbf{g} \\
    \boldsymbol{\tau}_{\text{total}} &= \sum_{i=0}^{3} \boldsymbol{\tau}_i
\end{align}
where $m_d = \sum_{i=0}^{3} m_i^{\text{drone}}$ is total drone mass.

\subsection{Drone Actuator Model}

\subsubsection{First-Order Thrust Dynamics}
Real quadrotor motors exhibit lag due to rotor inertia. We model this as a first-order system:
\begin{equation}
    \tau_m \dot{T} + T = T_c
    \label{eq:motor_continuous}
\end{equation}
where $T$ is actual thrust, $T_c$ is commanded thrust, and $\tau_m = \SI{0.2}{s}$ is the motor time constant.

Discretizing with forward Euler:
\begin{equation}
    T_{k+1} = T_k + \alpha(T_c - T_k), \quad \alpha = \min\left(\frac{\Delta t}{\tau_m}, 1\right)
    \label{eq:motor_discrete}
\end{equation}

\subsubsection{Thrust Saturation}
Each drone has maximum thrust $T_{\max} = \SI{250}{N}$. The thrust vector is saturated:
\begin{equation}
    \mathbf{T}_i \leftarrow \begin{cases}
        \mathbf{T}_i & \text{if } \|\mathbf{T}_i\| \leq T_{\max} \\
        T_{\max} \frac{\mathbf{T}_i}{\|\mathbf{T}_i\|} & \text{otherwise}
    \end{cases}
    \label{eq:thrust_saturation}
\end{equation}

\subsection{Noise and Disturbances}

\subsubsection{Drone Mass Variance}
Each drone mass is perturbed:
\begin{equation}
    m_i = m_{\text{nominal}} + \mathcal{N}(0, \sigma_m^2), \quad \sigma_m = 0.0001 \cdot m_{\text{nominal}}
    \label{eq:mass_variance}
\end{equation}

\subsubsection{Thrust Noise}
Thrust vectors are corrupted by Gaussian noise:
\begin{equation}
    \mathbf{T}_i \leftarrow \mathbf{T}_i + \mathcal{N}(\mathbf{0}, \sigma_T^2\mathbf{I}_3), \quad \sigma_T = 0.01 \|\mathbf{T}_i\|
    \label{eq:thrust_noise}
\end{equation}

\subsubsection{Attachment Point Error}
Attachment positions include fabrication errors:
\begin{equation}
    \mathbf{r}_i^{\text{actual}} = \mathbf{r}_i^{\text{nominal}} + \boldsymbol{\epsilon}_i, \quad \boldsymbol{\epsilon}_i \sim \mathcal{N}(\mathbf{0}, \sigma_a^2\mathbf{I}_3)
    \label{eq:attachment_error}
\end{equation}
with $\sigma_a = 0.0001 \cdot a$.

\subsection{Ground Contact Model}

A simple penetration-based contact model is employed. Bottom corners are:
\begin{equation}
    \mathbf{b}_j = \mathbf{p} + \mathbf{R}\left[\frac{a}{2}(\pm 1, \pm 1, -1)^T\right], \quad j \in \{0,1,2,3\}
\end{equation}

Contact occurs when $\min_j b_{j,z} < 0$. Upon contact:
\begin{enumerate}
    \item Project position: $\mathbf{p} \leftarrow \mathbf{p} - [0, 0, \min_j b_{j,z}]^T$
    \item Zero velocities: $\mathbf{v} \leftarrow \mathbf{0}$, $\boldsymbol{\omega} \leftarrow \mathbf{0}$
\end{enumerate}
This represents a perfectly inelastic collision.

% =============================================================================
\section{Hierarchical Control Architecture}
\label{sec:control}
% =============================================================================

\subsection{Control Overview}

The control system follows a hierarchical architecture:
\begin{enumerate}
    \item \textbf{Trajectory Planner}: Generates reference position, velocity, and feedforward acceleration
    \item \textbf{Position Controller (Outer Loop)}: PID control computing desired force
    \item \textbf{Attitude Controller (Inner Loop)}: Geometric control on SO(3) with thrust allocation
    \item \textbf{Thrust Saturation}: Enforces system-level constraints
\end{enumerate}

\subsection{Geometric Attitude Controller}

Following Lee et al. (2010), we implement geometric tracking control on SO(3).

\subsubsection{Attitude Error}
The attitude error is defined on SO(3):
\begin{equation}
    \mathbf{e}_R = \frac{1}{2}(\mathbf{R}_d^T\mathbf{R} - \mathbf{R}^T\mathbf{R}_d)^\vee
    \label{eq:attitude_error}
\end{equation}
where $(\cdot)^\vee : \mathfrak{so}(3) \to \mathbb{R}^3$ is the vee map (inverse of $[\cdot]_\times$).

\subsubsection{Control Law}
The desired torque is:
\begin{equation}
    \boldsymbol{\tau}_d = -K_R\mathbf{e}_R - K_\omega\mathbf{e}_\omega + \boldsymbol{\omega} \times \mathbf{I}\boldsymbol{\omega}
    \label{eq:attitude_control}
\end{equation}
where:
\begin{itemize}
    \item $K_R = 4 \cdot T_{\max}$ is the attitude error gain
    \item $K_\omega = 4 \cdot T_{\max}$ is the angular velocity error gain
    \item $\mathbf{e}_\omega = \boldsymbol{\omega} - \boldsymbol{\omega}_d$ is angular velocity error
    \item The feedforward term compensates gyroscopic effects
\end{itemize}

\subsubsection{Thrust Allocation}
We must map desired torque $\boldsymbol{\tau}_d$ to individual drone corrections $\delta\mathbf{F}_i$. The moment arm matrix is:
\begin{equation}
    \mathbf{A} = \begin{bmatrix} [\mathbf{r}_0]_\times \\ [\mathbf{r}_1]_\times \\ [\mathbf{r}_2]_\times \\ [\mathbf{r}_3]_\times \end{bmatrix} \in \mathbb{R}^{12 \times 3}
    \label{eq:moment_arm_matrix}
\end{equation}
where $\mathbf{r}_i = \mathbf{r}_i^{\text{world}} - \mathbf{p}$ are relative position vectors.

The least-squares solution is:
\begin{equation}
    \delta\mathbf{F} = \mathbf{A}^\dagger \boldsymbol{\tau}_d = (\mathbf{A}^T\mathbf{A})^{-1}\mathbf{A}^T\boldsymbol{\tau}_d
    \label{eq:thrust_allocation}
\end{equation}

\subsection{PID Position Controller}

\subsubsection{Control Law}
The position controller computes desired force via PID:
\begin{equation}
    \mathbf{F}_{\text{fb}} = \mathbf{K}_p\mathbf{e}_p + \mathbf{K}_d\mathbf{e}_v + \mathbf{K}_i\int \mathbf{e}_p \, dt
    \label{eq:pid}
\end{equation}
where:
\begin{align}
    \mathbf{e}_p &= \mathbf{p}_d - \mathbf{p} \\
    \mathbf{e}_v &= \mathbf{v}_d - \mathbf{v}
\end{align}

\subsubsection{Gain Tuning}
Gains are computed using Ziegler-Nichols tuning:
\begin{itemize}
    \item Ultimate gains: $K_{u,xy} = 0.08 T_{\max}$, $K_{u,z} = 0.12 T_{\max}$
    \item Ultimate periods: $T_{u,xy} = \SI{3.14}{s}$, $T_{u,z} = \SI{2.56}{s}$
\end{itemize}

PD control gains:
\begin{align}
    \mathbf{K}_p &= \text{diag}(0.8 K_{u,xy}, 0.8 K_{u,xy}, 0.8 K_{u,z}) \\
    \mathbf{K}_d &= \text{diag}(0.1 K_{u,xy}T_{u,xy}, 0.1 K_{u,xy}T_{u,xy}, 0.1 K_{u,z}T_{u,z})
\end{align}

Integral action is disabled ($\mathbf{K}_i = \mathbf{0}$) to avoid wind-up issues.

\subsection{Thrust Combination and Saturation}

\subsubsection{Force Decomposition}
Total force per drone combines feedforward and feedback terms:
\begin{equation}
    \mathbf{F}_i = \frac{1}{4}\mathbf{F}_{\text{ff}} + \frac{1}{4}\mathbf{F}_{\text{fb}} + \delta\mathbf{F}_i
    \label{eq:total_force}
\end{equation}
where $\mathbf{F}_{\text{ff}} = m\mathbf{a}_d$ is feedforward from trajectory.

\subsubsection{Control Authority Limits}
Available thrust is partitioned:
\begin{align}
    F_{\text{takeoff}} &= 0.75 \cdot 4T_{\max} = \SI{750}{N} \\
    F_{\text{transl}} &= 0.33 \cdot (4T_{\max} - F_{\text{takeoff}}) = \SI{82.5}{N} \\
    F_{\text{control}} &= 4T_{\max} - F_{\text{takeoff}} - F_{\text{transl}} = \SI{167.5}{N}
\end{align}

Control forces are saturated:
\begin{equation}
    \|\mathbf{F}_{\text{fb}} + \sum_i \delta\mathbf{F}_i\| \leq F_{\text{control}}
    \label{eq:control_saturation}
\end{equation}

% =============================================================================
\section{Trajectory Generation}
\label{sec:trajectory}
% =============================================================================

\subsection{Quintic Polynomial Trajectories}

\subsubsection{Boundary Conditions}
For a trajectory from $s_0 = 0$ to $s_f = d$ over duration $T$, we require:
\begin{align}
    s(0) = 0, \quad & s(T) = d \\
    \dot{s}(0) = 0, \quad & \dot{s}(T) = 0 \\
    \ddot{s}(0) = 0, \quad & \ddot{s}(T) = 0
\end{align}

\subsubsection{Polynomial Form}
The quintic polynomial is:
\begin{equation}
    s(t) = d(c_5t^5 + c_4t^4 + c_3t^3)
    \label{eq:quintic}
\end{equation}
where coefficients are derived from boundary conditions:
\begin{align}
    c_3 &= 10T^{-3} \\
    c_4 &= -15T^{-4} \\
    c_5 &= 6T^{-5}
\end{align}

Normalized time $\tau = t/T$ yields:
\begin{equation}
    s(\tau) = d(10\tau^3 - 15\tau^4 + 6\tau^5)
\end{equation}

\subsubsection{Derivatives}
Velocity and acceleration are:
\begin{align}
    \dot{s}(t) &= d \cdot T^{-1}(30\tau^2 - 60\tau^3 + 30\tau^4) \\
    \ddot{s}(t) &= d \cdot T^{-2}(60\tau - 180\tau^2 + 120\tau^3)
\end{align}

\subsection{Multi-Phase Trajectory Planning}

The mission consists of seven phases:

\begin{table}[h]
\centering
\caption{Mission phases and durations}
\begin{tabular}{clr}
\toprule
Phase & Description & Duration \\
\midrule
0 & Takeoff (vertical to \SI{5}{m}) & $T_0$ \\
1 & Hover at altitude & \SI{2}{s} \\
2 & Translate East (\SI{10}{m}) & $T_2$ \\
3 & Hover at waypoint & \SI{2}{s} \\
4 & Translate North (\SI{10}{m}) & $T_4$ \\
5 & Final hover & \SI{2}{s} \\
6 & Landing (controlled descent) & $T_6$ \\
\bottomrule
\end{tabular}
\label{tab:phases}
\end{table}

\subsubsection{Trajectory Duration Computation}
Phase durations are computed from acceleration limits:
\begin{equation}
    T_i = \sqrt{\frac{10\sqrt{3}d_i}{3a_{\max,i}}}
    \label{eq:phase_duration}
\end{equation}
where $d_i$ is travel distance and $a_{\max,i}$ is maximum acceleration for that phase.

\subsubsection{Reference Signals}
At each timestep, the planner provides:
\begin{itemize}
    \item $\mathbf{p}_d(t)$: Reference position
    \item $\mathbf{v}_d(t)$: Reference velocity
    \item $\mathbf{a}_d(t)$: Feedforward acceleration
\end{itemize}

These are computed by evaluating quintic polynomials for the active phase.

% =============================================================================
\section{Implementation}
\label{sec:implementation}
% =============================================================================

\subsection{Software Architecture}

The simulation is implemented in Python 3.13 with the following structure:

\begin{verbatim}
csd/
├── main.py           # Main simulation loop and payload dynamics
├── drone.py          # Drone actuator model
├── entity.py         # Base class for simulated entities
├── config.py         # Configuration parameters
├── clock.py          # Simulation time management
├── controller/
│   ├── thrust.py     # Main controller coordinator
│   ├── position.py   # PID position controller
│   ├── attitude.py   # Geometric attitude controller
│   ├── plan.py       # Trajectory planner
│   └── quintic.py    # Quintic polynomial functions
├── receiver.py       # Telemetry receiver
└── viz/              # Visualization modules
\end{verbatim}

\subsection{Key Dependencies}

\begin{itemize}
    \item \texttt{numpy 2.3.4}: Numerical computations
    \item \texttt{scipy 1.16.3}: Rotation representations and integration
    \item \texttt{matplotlib 3.10.7}: Plotting and visualization
    \item \texttt{pandas 2.3.3}: Data analysis
    \item \texttt{pydantic 2.12.4}: Configuration validation
\end{itemize}

\subsection{Numerical Integration}

\subsubsection{Time Step}
The simulation uses fixed time step $\Delta t = \SI{0.001}{s}$ (1 kHz update rate).

\subsubsection{Integration Scheme}
Forward Euler integration is employed:
\begin{align}
    \mathbf{v}_{k+1} &= \mathbf{v}_k + \mathbf{a}_k \Delta t \\
    \mathbf{p}_{k+1} &= \mathbf{p}_k + \mathbf{v}_k \Delta t \\
    \boldsymbol{\omega}_{k+1} &= \boldsymbol{\omega}_k + \boldsymbol{\alpha}_k \Delta t \\
    \mathbf{R}_{k+1} &= \mathbf{R}_k \exp([\boldsymbol{\omega}_k \Delta t]_\times)
\end{align}

\textbf{Note:} Forward Euler is first-order accurate and does not conserve energy. For production systems, consider Runge-Kutta or symplectic integrators.

\subsection{Telemetry System}

\subsubsection{Data Transmission}
Each entity transmits telemetry at each timestep:
\begin{itemize}
    \item Payload: position, velocity, acceleration, orientation, angular velocity
    \item Drones: position, commanded thrust, actual thrust
\end{itemize}

\subsubsection{Data Format}
Telemetry is serialized to JSON and transmitted via message queues for real-time monitoring and post-processing.

\subsection{Configuration Management}

All parameters are centralized in \texttt{config.py}:
\begin{itemize}
    \item \textbf{Feature switches}: Enable/disable noise sources and controllers
    \item \textbf{Physical parameters}: Masses, dimensions, gravity
    \item \textbf{Control gains}: PID and geometric controller parameters
    \item \textbf{Trajectory waypoints}: Mission-specific coordinates
    \item \textbf{Noise parameters}: Variance and scaling factors
\end{itemize}

This design enables parameter sweeps and ablation studies without code modification.

% =============================================================================
\section{Experimental Results}
\label{sec:experiments}
% =============================================================================

% STUB SECTION - TO BE COMPLETED

\subsection{Experimental Setup}

\subsubsection{Simulation Parameters}
All experiments use the following baseline configuration:
\begin{itemize}
    \item Simulation duration: \SI{30}{s}
    \item Time step: \SI{0.001}{s}
    \item Payload mass: \SI{10}{kg}
    \item Drone max thrust: \SI{250}{N} each
    \item Target waypoints: $(0, 0, 5)$, $(10, 0, 5)$, $(10, 10, 5)$ meters
\end{itemize}

\subsubsection{Performance Metrics}
We evaluate performance using:
\begin{itemize}
    \item Position tracking error: $e_p(t) = \|\mathbf{p}(t) - \mathbf{p}_d(t)\|$
    \item Velocity magnitude: $\|\mathbf{v}(t)\|$
    \item Attitude error: Euler angle deviations from level flight
    \item Settling time: Time to reach $e_p < \SI{0.01}{m}$ after maneuver
    \item Overshoot: Maximum position error during transition
\end{itemize}

\subsection{Baseline Performance}

% TODO: Add results from nominal simulation run
% - Trajectory tracking plots (x, y, z vs time)
% - Velocity profiles
% - Attitude evolution (roll, pitch, yaw)
% - Control effort (thrust magnitudes)
% - Phase transition behavior

\subsubsection{Trajectory Tracking}
% STUB: Insert Figure - Position vs. time for all three axes
% STUB: Insert Table - Tracking error statistics (mean, max, RMS)

\subsubsection{Attitude Regulation}
% STUB: Insert Figure - Euler angles vs. time
% STUB: Insert Table - Maximum attitude deviations

\subsubsection{Control Effort}
% STUB: Insert Figure - Individual drone thrust magnitudes
% STUB: Insert Figure - Total system thrust vs. available thrust

\subsection{Noise Sensitivity Analysis}

% TODO: Sweep noise parameters and measure degradation

\subsubsection{Mass Variance Effects}
% STUB: Vary DRONE_MASS_VARIANCE from 0 to 0.01
% STUB: Plot tracking error vs. variance
% STUB: Analyze asymmetric loading effects

\subsubsection{Thrust Noise Effects}
% STUB: Vary DRONE_THRUST_NOISE from 0 to 0.05
% STUB: Measure high-frequency oscillations
% STUB: Compare filtered vs. raw thrust signals

\subsubsection{Attachment Error Effects}
% STUB: Vary ATTACHMENT_ERROR_SCALE from 0 to 0.001
% STUB: Analyze induced torque disturbances
% STUB: Measure attitude controller compensation

\subsection{Controller Gain Tuning}

% TODO: Ablation studies on control gains

\subsubsection{Position Controller Gains}
% STUB: Sweep KP, KD values
% STUB: Analyze stability margins
% STUB: Compare against Ziegler-Nichols predictions

\subsubsection{Attitude Controller Gains}
% STUB: Sweep KR, KW values
% STUB: Measure attitude response time
% STUB: Identify optimal damping ratio

\subsection{Ablation Studies}

\subsubsection{Controller Disable Experiments}
% STUB: Run with ENABLE_ATTITUDE_CONTROLLER = False
% STUB: Run with ENABLE_POSITION_CONTROLLER = False
% STUB: Quantify contribution of each controller

\subsubsection{Noise Disable Experiments}
% STUB: Compare performance with all noise disabled
% STUB: Establish upper bound on achievable performance
% STUB: Measure noise-induced degradation

\subsection{Payload Mass Variation}

% TODO: Vary CUBE_MASS and measure performance

\subsubsection{Light Payload}
% STUB: CUBE_MASS = 5 kg
% STUB: Measure improved agility
% STUB: Analyze reduced settling time

\subsubsection{Heavy Payload}
% STUB: CUBE_MASS = 20 kg
% STUB: Measure thrust saturation effects
% STUB: Analyze trajectory tracking degradation

\subsection{Trajectory Complexity}

% TODO: Test with more aggressive trajectories

\subsubsection{Increased Translation Distance}
% STUB: EAST = NORTH = 20 m
% STUB: Measure sustained velocity and acceleration
% STUB: Analyze control effort distribution

\subsubsection{Reduced Hover Times}
% STUB: HOVER_TIMES = [0.5, 0.5, 0.5]
% STUB: Test rapid maneuver transitions
% STUB: Measure transient response

\subsection{Computational Performance}

\subsubsection{Profiling Results}
% STUB: Present cProfile output from ENABLE_PROFILING
% STUB: Identify computational bottlenecks
% STUB: Report real-time factor achieved

\subsubsection{Scalability}
% STUB: Measure performance vs. time step size
% STUB: Analyze stability for larger dt
% STUB: Compare accuracy vs. computational cost

\subsection{Comparison with Literature}

% STUB: Compare results against published cooperative transport work
% STUB: Benchmark against Lee 2010 tracking performance
% STUB: Discuss advantages/limitations of approach

% =============================================================================
\section{Discussion}
% =============================================================================

\subsection{Key Findings}

% STUB: Summarize experimental results
% - Hierarchical control successfully stabilizes system
% - Geometric attitude control avoids singularities
% - Quintic trajectories provide smooth motion
% - System robust to moderate noise levels

\subsection{Limitations}

\subsubsection{Modeling Assumptions}
\begin{itemize}
    \item Perfect knowledge of payload inertia
    \item Rigid attachment assumption (no cable dynamics)
    \item Simple ground contact model
    \item No aerodynamic effects or rotor interactions
\end{itemize}

\subsubsection{Control Limitations}
\begin{itemize}
    \item Manual gain tuning (no adaptive control)
    \item No state estimation (assumes perfect measurements)
    \item Simple saturation handling (could use anti-windup schemes)
    \item No collision avoidance between drones
\end{itemize}

\subsection{Future Work}

\subsubsection{Enhanced Dynamics}
\begin{itemize}
    \item Cable suspension model for flexible attachment
    \item Aerodynamic drag and rotor interference
    \item Deformable payload dynamics
    \item Multi-payload formations
\end{itemize}

\subsubsection{Advanced Control}
\begin{itemize}
    \item Model predictive control (MPC) for constraint handling
    \item Adaptive control for parameter uncertainty
    \item Observer-based state estimation (EKF/UKF)
    \item Optimal trajectory planning
\end{itemize}

\subsubsection{Hardware Validation}
\begin{itemize}
    \item Hardware-in-the-loop (HIL) testing
    \item Flight experiments with real quadrotors
    \item Vision-based state estimation
    \item Distributed control implementation
\end{itemize}

% =============================================================================
\section{Conclusion}
\label{sec:conclusion}
% =============================================================================

This work presented a comprehensive simulation framework for cooperative quadrotor payload transport with hierarchical control. The system integrates high-fidelity 6-DOF rigid body dynamics, geometric attitude control on SO(3), and quintic polynomial trajectory generation.

Key contributions include:
\begin{enumerate}
    \item Complete dynamics simulation with realistic actuator models and noise
    \item Provably stable geometric attitude controller avoiding singularities
    \item Systematic trajectory generation ensuring smooth motion profiles
    \item Configurable architecture enabling extensive parametric studies
\end{enumerate}

Experimental results (to be completed) demonstrate robust performance under realistic operating conditions. The framework provides a valuable tool for algorithm development, controller tuning, and mission planning for cooperative aerial manipulation systems.

The open-source implementation facilitates reproducibility and serves as an educational resource for understanding multibody dynamics, Lie group control, and hierarchical control architectures.

% =============================================================================
\section*{Acknowledgments}
% =============================================================================

% STUB: Add acknowledgments as appropriate

% =============================================================================
% REFERENCES
% =============================================================================

\begin{thebibliography}{99}

\bibitem{lee2010}
Lee, T., Leok, M., \& McClamroch, N. H. (2010).
\textit{Geometric tracking control of a quadrotor UAV on SE(3)}.
IEEE Conference on Decision and Control (CDC), 5420--5425.

\bibitem{ziegler1942}
Ziegler, J. G., \& Nichols, N. B. (1942).
\textit{Optimum settings for automatic controllers}.
Transactions of the ASME, 64(11), 759--768.

\bibitem{mellinger2011}
Mellinger, D., \& Kumar, V. (2011).
\textit{Minimum snap trajectory generation and control for quadrotors}.
IEEE International Conference on Robotics and Automation (ICRA), 2520--2525.

\bibitem{bullo2004}
Bullo, F., \& Murray, R. M. (2004).
\textit{Geometric control of mechanical systems: Modeling, analysis, and design for simple mechanical control systems}.
Springer Science \& Business Media.

\bibitem{michael2011}
Michael, N., Fink, J., \& Kumar, V. (2011).
\textit{Cooperative manipulation and transportation with aerial robots}.
Autonomous Robots, 30(1), 73--86.

\bibitem{sreenath2013}
Sreenath, K., Lee, T., \& Kumar, V. (2013).
\textit{Geometric control and differential flatness of a quadrotor UAV with a cable-suspended load}.
IEEE Conference on Decision and Control (CDC), 2269--2274.

\bibitem{cruz2015}
Cruz, P. J., \& Fierro, R. (2015).
\textit{Cable-suspended load lifting by a quadrotor UAV: Hybrid model, trajectory generation, and control}.
Autonomous Robots, 41(8), 1629--1643.

\bibitem{jiang2017}
Jiang, Q., \& Kumar, V. (2017).
\textit{The inverse kinematics of cooperative transport with multiple aerial robots}.
IEEE Transactions on Robotics, 29(1), 136--145.

\bibitem{jackson2020}
Jackson, B. E., Howell, T. A., Shah, K., Schwager, M., \& Manchester, Z. (2020).
\textit{Scalable cooperative transport of cable-suspended loads with UAVs using distributed trajectory optimization}.
IEEE Robotics and Automation Letters, 5(2), 3368--3374.

\bibitem{pounds2014}
Pounds, P. E., Bersak, D. R., \& Dollar, A. M. (2014).
\textit{Stability of small-scale UAV helicopters and quadrotors with added payload mass under PID control}.
Autonomous Robots, 33(1), 129--142.

\end{thebibliography}

% =============================================================================
% APPENDICES
% =============================================================================

\appendix

\section{Mathematical Derivations}
\label{app:derivations}

\subsection{Quintic Polynomial Coefficient Derivation}

Starting from the general quintic:
\begin{equation}
    s(t) = a_0 + a_1 t + a_2 t^2 + a_3 t^3 + a_4 t^4 + a_5 t^5
\end{equation}

Boundary conditions at $t=0$:
\begin{align}
    s(0) = 0 &\Rightarrow a_0 = 0 \\
    \dot{s}(0) = 0 &\Rightarrow a_1 = 0 \\
    \ddot{s}(0) = 0 &\Rightarrow a_2 = 0
\end{align}

Boundary conditions at $t=T$:
\begin{align}
    s(T) = d &\Rightarrow a_3 T^3 + a_4 T^4 + a_5 T^5 = d \\
    \dot{s}(T) = 0 &\Rightarrow 3a_3 T^2 + 4a_4 T^3 + 5a_5 T^4 = 0 \\
    \ddot{s}(T) = 0 &\Rightarrow 6a_3 T + 12a_4 T^2 + 20a_5 T^3 = 0
\end{align}

Solving this system yields:
\begin{align}
    a_3 &= 10dT^{-3} \\
    a_4 &= -15dT^{-4} \\
    a_5 &= 6dT^{-5}
\end{align}

\subsection{Moment of Inertia for Uniform Cube}

For a uniform cube with density $\rho$, side length $a$, and mass $m = \rho a^3$:
\begin{equation}
    I_{xx} = \int\int\int (y^2 + z^2) \rho \, dx\,dy\,dz
\end{equation}

By symmetry, $I_{xx} = I_{yy} = I_{zz}$ and off-diagonal terms vanish. Evaluating:
\begin{equation}
    I_{xx} = \rho \int_{-a/2}^{a/2} \int_{-a/2}^{a/2} \int_{-a/2}^{a/2} (y^2 + z^2) \, dx\,dy\,dz = \frac{ma^2}{6}
\end{equation}

Therefore:
\begin{equation}
    \mathbf{I} = \frac{ma^2}{6} \begin{bmatrix} 1 & 0 & 0 \\ 0 & 1 & 0 \\ 0 & 0 & 1 \end{bmatrix}
\end{equation}

\section{Implementation Details}
\label{app:implementation}

\subsection{Rodrigues Formula Implementation}

The matrix exponential $\exp([\boldsymbol{\omega}]_\times)$ is computed via:
\begin{equation}
    \exp([\boldsymbol{\omega}]_\times) = \mathbf{I} + \frac{\sin\theta}{\theta}[\boldsymbol{\omega}]_\times + \frac{1-\cos\theta}{\theta^2}[\boldsymbol{\omega}]_\times^2
\end{equation}
where $\theta = \|\boldsymbol{\omega}\|$.

This is implemented in \texttt{scipy.spatial.transform.Rotation.from\_rotvec()}.

\subsection{Least-Squares Thrust Allocation Algorithm}

\begin{algorithm}
\caption{Least-Squares Thrust Allocation}
\begin{algorithmic}[1]
\Require Desired torque $\boldsymbol{\tau}_d$, attachment points $\{\mathbf{r}_i\}$
\Ensure Thrust corrections $\{\delta\mathbf{F}_i\}$
\State Construct moment arm matrix: $\mathbf{A}_{3i:3i+3,:} \gets [\mathbf{r}_i]_\times$ for $i=0,\ldots,3$
\State Compute pseudoinverse: $\mathbf{A}^\dagger \gets (\mathbf{A}^T\mathbf{A})^{-1}\mathbf{A}^T$
\State Allocate forces: $\delta\mathbf{F} \gets \mathbf{A}^\dagger \boldsymbol{\tau}_d$
\State Reshape to per-drone forces: $\delta\mathbf{F}_i \gets \delta\mathbf{F}[3i:3i+3]$
\State \Return $\{\delta\mathbf{F}_i\}$
\end{algorithmic}
\end{algorithm}

\subsection{Configuration Parameters}

\begin{table}[h]
\centering
\caption{Complete simulation parameter listing}
\begin{tabular}{lll}
\toprule
Parameter & Symbol & Value \\
\midrule
Time step & $\Delta t$ & \SI{0.001}{s} \\
Simulation duration & $T_{\text{sim}}$ & \SI{30}{s} \\
Payload mass & $m_c$ & \SI{10}{kg} \\
Payload size & $a$ & \SI{1}{m} \\
Drone mass & $m_d$ & \SI{1}{kg} \\
Drone max thrust & $T_{\max}$ & \SI{250}{N} \\
Motor time constant & $\tau_m$ & \SI{0.2}{s} \\
Gravity & $g$ & \SI{9.80665}{m/s^2} \\
Attitude gain & $K_R$ & $4T_{\max}$ \\
Angular velocity gain & $K_\omega$ & $4T_{\max}$ \\
Position gain (xy) & $K_{p,xy}$ & $0.8 \times 0.08 T_{\max}$ \\
Position gain (z) & $K_{p,z}$ & $0.8 \times 0.12 T_{\max}$ \\
Derivative gain (xy) & $K_{d,xy}$ & $0.1 \times 0.08 T_{\max} \times 3.14$ \\
Derivative gain (z) & $K_{d,z}$ & $0.1 \times 0.12 T_{\max} \times 2.56$ \\
Mass variance & $\sigma_m$ & $0.0001 m_d$ \\
Thrust noise & $\sigma_T$ & $0.01 \|T\|$ \\
Attachment error & $\sigma_a$ & $0.0001 a$ \\
\bottomrule
\end{tabular}
\label{tab:parameters}
\end{table}

\section{Telemetry Data Format}
\label{app:telemetry}

\subsection{Payload Telemetry Schema}

\begin{lstlisting}[language=json]
{
  "topic": "cube",
  "time": 1.234,
  "pos": [x, y, z],
  "vel": [vx, vy, vz],
  "acc": [ax, ay, az],
  "rot": [[r11, r12, r13],
          [r21, r22, r23],
          [r31, r32, r33]],
  "ang_vel": [wx, wy, wz],
  "ang_acc": [alpha_x, alpha_y, alpha_z]
}
\end{lstlisting}

\subsection{Drone Telemetry Schema}

\begin{lstlisting}[language=json]
{
  "topic": "drone_0",
  "time": 1.234,
  "pos": [x, y, z],
  "command": [Tcx, Tcy, Tcz],
  "thrust": [Tx, Ty, Tz]
}
\end{lstlisting}

\end{document}
